\documentclass[11pt,a4paper]{article}
\usepackage{apnotes}
\setsunday{1}
\begin{document}

\notestitle{Data \& Machine Learning Foundations}{1--5}{Think like a data scientist: frame a business problem, clean data, build and explain a model, track experiments, and ship project \textbf{P1 --- Churn Intelligence} to GitHub by 18:00.}

\begin{tabularx}{\linewidth}{@{}c c X l@{}}
\toprule
\textbf{Hour} & \textbf{Mode} & \textbf{Focus} & \textbf{Dataset}\\
\midrule
1 & \teachtag & Scientific Python and exploratory data analysis (EDA) & \ds{UCI Adult Income}\\
2 & \teachtag & ML pipelines, cross-validation, gradient boosting & \ds{California Housing}\\
3 & \buildtag & Churn project: problem framing and a baseline & \ds{IBM Telco Churn}\\
4 & \buildtag & Class imbalance, Optuna tuning, SHAP explanations & \ds{IBM Telco Churn}\\
5 & \buildtag & MLflow tracking, repository, stakeholder summary & \ds{Telco + MLflow runs}\\
\bottomrule
\end{tabularx}

\medskip
Companion notebook: \texttt{notebooks/S01\_Data\_ML\_Foundations.ipynb}. Every number in a green box below was produced by running it.

%% ------------------------------------------------------------------
\hourhead{1}{TEACH}{Scientific Python \& Exploratory Data Analysis}{UCI Adult Income (48,842 rows)}{NumPy, Pandas, Matplotlib, data cleaning, EDA reasoning}

\subsection{Why vectorisation matters}
A Python loop handles one number at a time. NumPy applies the same operation to a whole array in compiled code.
\begin{lstlisting}
x = np.random.rand(1_000_000)
s_loop = sum(v * v for v in x)      # slow: interpreted loop
s_vec  = np.dot(x, x)               # fast: one call into optimised C
\end{lstlisting}
\begin{resultbox}
Python loop 135.2\,ms versus NumPy 0.9\,ms --- about \textbf{154$\times$ faster}, same answer.
\end{resultbox}

\subsection{A data-quality checklist}
Before any model, answer five questions about every dataset:
\begin{enumerate}
  \item \textbf{Shape and types} --- how many rows and columns; which are numeric, which categorical?
  \item \textbf{Missing values} --- how many, in which columns, and \emph{why} are they missing?
  \item \textbf{Duplicates} --- exact repeated rows inflate evaluation scores.
  \item \textbf{Target balance} --- what share of rows are positive?
  \item \textbf{Leakage} --- does any column contain information you would not have at prediction time?
\end{enumerate}
\begin{resultbox}
Adult: 48,842 rows. Missing (``?''): \texttt{workclass} 5.73\%, \texttt{occupation} 5.75\%, \texttt{native\_country} 1.75\%. 52 duplicate rows removed. After cleaning, 23.9\% of people earn more than \$50K.
\end{resultbox}

\begin{keyidea}[Missing is information]
Filling \texttt{occupation} with the mode hides a pattern. Making ``Unknown'' an explicit category keeps the rows \emph{and} lets the model learn whether missingness itself is predictive. Write every cleaning decision in your README.
\end{keyidea}

\figw[0.98]{s1_adult_eda.png}{EDA on Adult Income: share above \$50K rises with years of education; the higher-income group is older.}

\begin{resultbox}
Share earning more than \$50K: \textbf{30.4\%} of men versus \textbf{10.9\%} of women. This gap is a question to investigate, not just a feature --- we return to it in the Sunday 6 fairness audit.
\end{resultbox}

%% ------------------------------------------------------------------
\hourhead{2}{TEACH}{Supervised ML pipelines, cross-validation, gradient boosting}{California Housing (20,640 districts)}{scikit-learn Pipelines, k-fold CV, XGBoost, feature engineering}

\subsection{The golden rule: never touch the test set until the end}
Split once into training (80\%) and test (20\%). Choose models using $k$-fold cross-validation on the training part only:
\[
\mathrm{CV\ RMSE}=\frac{1}{k}\sum_{i=1}^{k}\sqrt{\frac{1}{|F_i|}\sum_{j\in F_i}(y_j-\hat y_j)^2}
\]
where $F_i$ is the $i$-th fold, scored by a model trained on the other $k-1$ folds.

\subsection{Pipelines prevent leakage}
A \texttt{Pipeline} learns imputation and scaling \emph{inside each fold}, so statistics from validation rows never leak into training.
\begin{lstlisting}
prep = ColumnTransformer([
  ("num", Pipeline([("impute", SimpleImputer(strategy="median")),
                    ("scale", StandardScaler())]), num_cols),
  ("cat", OneHotEncoder(handle_unknown="ignore"), ["ocean_proximity"])])
pipe = Pipeline([("prep", prep), ("model", XGBRegressor(n_estimators=400))])
\end{lstlisting}

\textbf{Feature engineering.} Ratios such as rooms per household, bedrooms per room and people per household carry more signal than raw totals.

\begin{resultbox}[From our notebook run: 5-fold CV on the training set]
\begin{tabular}{lrr}
\toprule
Model & CV RMSE (\$) & std\\
\midrule
XGBoost & 45,339 & 967\\
Random forest & 50,258 & 1,038\\
Linear regression & 67,859 & 1,312\\
\bottomrule
\end{tabular}\\[4pt]
Final XGBoost on the untouched test set: \textbf{RMSE \$44,355, $R^2=0.850$}.
\end{resultbox}

\figw[0.48]{s1_housing_pred.png}{Predicted versus actual house values. The horizontal band at \$500,001 is a cap in the data, not a model failure.}

\begin{pitfall}[Capped targets]
California Housing truncates expensive districts at \$500,001. Spot such artefacts during EDA and say how you handled them (keep, drop, or model separately).
\end{pitfall}

%% ------------------------------------------------------------------
\hourhead{3}{BUILD}{Churn project: problem framing and a baseline}{IBM Telco Customer Churn (7,043 customers)}{Problem framing, baseline thinking, F1 / ROC-AUC}

\subsection{Frame the problem before you model it}
\begin{tabularx}{\linewidth}{@{}l X@{}}
\textbf{Business question} & Which customers are likely to cancel next month, so the retention team can call them first?\\
\textbf{Unit of prediction} & One customer.\\
\textbf{Target} & \texttt{Churn = Yes}.\\
\textbf{Metrics} & ROC-AUC (ranking quality) and F1 for the churn class.\\
\end{tabularx}

\begin{defbox}[Metrics for imbalanced classes]
Precision $P=\frac{TP}{TP+FP}$, recall $R=\frac{TP}{TP+FN}$, and $F_1=\frac{2PR}{P+R}$.
ROC-AUC is the probability that a random churner gets a higher score than a random non-churner.
\end{defbox}

\begin{resultbox}
Churn rate 26.5\%. By contract: month-to-month 42.7\%, one-year 11.3\%, two-year 2.8\%.\\
Majority-class dummy: ROC-AUC 0.500, F1 0.000 --- yet about 73\% accuracy.\\
Logistic regression: ROC-AUC 0.846, F1 0.607.
\end{resultbox}

\begin{pitfall}[Accuracy hides failure]
A model that never predicts churn is 73\% accurate and completely useless. On imbalanced problems, always report a class-aware metric and compare against a dummy baseline.
\end{pitfall}

%% ------------------------------------------------------------------
\hourhead{4}{BUILD}{Improve the model: imbalance, Optuna, SHAP}{IBM Telco Customer Churn}{Optuna, class weighting, SHAP explainability}

\textbf{Class weighting.} Instead of oversampling, weight positive examples by $\frac{\#\text{negatives}}{\#\text{positives}}=2.77$ (\texttt{scale\_pos\_weight}).

\textbf{Optuna} searches hyperparameters with a Bayesian sampler (TPE): it proposes settings, scores them with cross-validation, and focuses on promising regions.
\begin{lstlisting}
def objective(trial):
    params = dict(max_depth=trial.suggest_int("max_depth", 2, 6),
                  learning_rate=trial.suggest_float("learning_rate", 0.01, 0.2, log=True), ...)
    return cross_val_score(make_pipe(params), X_tr, y_tr, cv=4, scoring="roc_auc").mean()
study = optuna.create_study(direction="maximize"); study.optimize(objective, n_trials=25)
\end{lstlisting}

\begin{resultbox}[Leaderboard on the held-out 25\% test set]
\begin{tabular}{lrr}
\toprule
Model & ROC-AUC & F1 (churn)\\
\midrule
XGBoost + class weight + Optuna & \textbf{0.847} & \textbf{0.632}\\
Logistic regression + class weight & 0.846 & 0.627\\
Logistic regression & 0.846 & 0.607\\
Majority-class dummy & 0.500 & 0.000\\
\bottomrule
\end{tabular}
\end{resultbox}

\begin{keyidea}[Honest reading of the leaderboard]
Twenty-five rounds of tuning bought only 0.001 ROC-AUC over a plain logistic regression. Class weighting helped F1 more than tuning did. When a simple model is this close, prefer it: it is faster, easier to explain and easier to monitor.
\end{keyidea}

\subsection{SHAP: explaining the model to the business}
SHAP splits each prediction into contributions from each feature, so that $\hat f(x)=\phi_0+\sum_j\phi_j(x)$.
\figw[0.62]{s1_shap.png}{Top churn drivers by mean $|\text{SHAP}|$: month-to-month contract, short tenure, fibre-optic internet, no online security.}

\begin{pitfall}[SHAP is not causation]
SHAP shows what the \emph{model} relies on, not what \emph{causes} churn. Use it to suggest actions (a loyalty offer before month 12), then test those actions with an experiment.
\end{pitfall}

%% ------------------------------------------------------------------
\hourhead{5}{BUILD}{Track experiments, publish, present}{Telco + MLflow runs}{MLflow, Git/GitHub workflow, technical storytelling}

\begin{lstlisting}
mlflow.set_tracking_uri("sqlite:///mlflow.db")
mlflow.set_experiment("P1-churn-intelligence")
with mlflow.start_run(run_name="XGBoost + class weight + Optuna"):
    mlflow.log_param("model", "xgboost"); mlflow.log_metrics({"roc_auc": 0.847, "f1_churn": 0.632})
\end{lstlisting}
Open the dashboard with \texttt{mlflow ui --backend-store-uri sqlite:///mlflow.db}.

\textbf{One-page summary for a manager.} Question; data; best model and its held-out score; top three drivers; one recommendation; limits (one company, one snapshot, SHAP is not causation).

\begin{exercise}
\begin{enumerate}
  \item Add \texttt{tenure} buckets (0--12, 13--24, 25+ months) as a feature. Does F1 change?
  \item Choose the probability threshold that maximises F1 on a validation split instead of using 0.5.
  \item Estimate the business value: if a call costs Rs.\,50 and a saved customer is worth Rs.\,3,000, which threshold maximises profit?
\end{enumerate}
\end{exercise}

\begin{submitbox}
Deliverable \textbf{P1 --- Churn Intelligence}: notebook, MLflow logs, SHAP plot and one-page summary in \texttt{P1-churn-intelligence/}.\\
\texttt{git push origin main} $\rightarrow$ \texttt{git tag -a sunday-1 -m "Sunday 1 submission"} $\rightarrow$ \texttt{git push origin sunday-1} $\rightarrow$ create a GitHub Release $\rightarrow$ submit the link.
\end{submitbox}

\section*{Checklist before you leave}
\begin{itemize}
  \item[$\square$] Cleaning decisions written in the README
  \item[$\square$] Dummy baseline reported next to every model
  \item[$\square$] Test set used once, at the end
  \item[$\square$] SHAP drivers translated into one business action
  \item[$\square$] Tag \texttt{sunday-1} pushed before 18:00
\end{itemize}

\end{document}
